\documentclass{article}
\usepackage[utf8]{inputenc}
\usepackage{amsmath}
\usepackage{geometry}
\geometry{margin=2cm}
\begin{document}

\section*{Questão 116 — ENEM 2024 — Ciências da Natureza}

\subsection*{Interpretação e Estratégia}

\textbf{Introdução:} Esta questão propõe a análise de um cladograma que apresenta três grupos distintos dentro dos mamíferos: Monotremata, Marsupialia e Placentalia. O foco da questão é identificar as mudanças evolutivas nos padrões de reprodução e nutrição dos embriões e descendentes que ocorrem entre esses grupos.

\textbf{Termos-chave:}
\begin{itemize}
  \item \textbf{Cladograma:} diagrama que representa relações evolutivas.
  \item \textbf{Evolução:} modificação genética ao longo do tempo.
  \item \textbf{Nutrição embrionária:} forma de suprir nutrientes durante o desenvolvimento.
\end{itemize}

\textbf{O que a questão pede:} identificar a perda e ganho evolutivos relacionados à reprodução e nutrição entre os grupos do cladograma.

\textbf{Estratégia de resolução:} interpretar corretamente o cladograma e aplicar conhecimento sobre evolução dos mamíferos, especialmente redução do vitelo (vesícula vitelínica) e surgimento da lactação.

\subsection*{Conteúdo e Mini Revisão}

\textbf{Tabela de conceitos:}

\begin{tabular}{p{4cm}p{6cm}p{6cm}}
\hline
\textbf{Conceito} & \textbf{Ideia-chave} & \textbf{Aplicação na questão} \\
\hline
Cladograma & Relações evolutivas entre grupos & Identificar grupos mamíferos e características reprodutivas \\
Vesícula vitelínica & Armazena reservas nutricionais no ovo & Redução implica maior dependência da lactação \\
Lactação & Nutrição dos filhotes via leite materno & Inovação evolutiva nos mamíferos explorada no cladograma \\
\hline
\end{tabular}

\bigskip

\textbf{Fundamentação teórica:}
\begin{itemize}
  \item Mamíferos apresentam reprodução sexuada com fecundação interna e geralmente desenvolvimento interno dos embriões.
  \item Grupos apresentados: \emph{Monotremata} (mamíferos ovíparos com vesícula vitelínica maior), \emph{Marsupialia} (desenvolvimento incompleto na placenta, grande mudança no cuidado pós-natal), \emph{Placentalia} (nutrição intrauterina eficiente e lactação robusta).
  \item Evolutivamente, houve redução do vitelo no ovo compensada pelo desenvolvimento da lactação.
  \item Essa transição representa perdas (vesícula vitelínica) e ganhos (lactação) adaptativos importantes para nutrição e sobrevivência.
\end{itemize}

\subsection*{Resolução e "Pulo do Gato"}

\begin{enumerate}
  \item Examine o cladograma: identifique os grupos Monotremata, Marsupialia e Placentalia.
  \item Reconheça que todos os grupos têm reprodução sexuada e fecundação interna, eliminando opções contrárias.
  \item Observe a redução da vesícula vitelínica e o desenvolvimento da lactação como principais mudanças reprodutivas.
  \item Selecione a alternativa que aborda essas características — alternativa \textbf{E}.
\end{enumerate}

\bigskip

\textbf{Pulo do Gato:} \emph{Foque na relação entre a redução do vitelo e o desenvolvimento da lactação para identificar o avanço evolutivo nos grupos de mamíferos.}

\subsection*{Armadilhas e Distratores}

\begin{itemize}
  \item \textbf{A:} interpreta aumento de descendentes como avanço principal, mas foco é nutrição.
  \item \textbf{B:} sugere mudança na fecundação que já é interna neste grupo.
  \item \textbf{C:} supõe reprodução se tornou exclusivamente sexuada, incorreto pois já era.
  \item \textbf{D:} associa evolução à migração aquática-terrestre do desenvolvimento embrionário, não pertinente.
\end{itemize}

\subsection*{Código da Questão}

CN\_REPRODUCAO\_CIENCIASDA NATUREZA\_001

\bigskip

\section*{Referência do Cladograma}

\textit{DIXSON, A. F. Mammalian Sexuality: The Act of Mating and the Evolution of Reproduction.} Disponível em: www.cambridge.org. Acesso em: 2 jul. 2024 (adaptado).

\bigskip

% Imagens do enunciado não reproduzidas em LaTeX pela política solicitada, mas ferramentas visuais do cladograma são essenciais para apoio na resolução.

\end{document}